% Generated draft: requires mathematical and source review.
\documentclass[11pt]{article}
\usepackage[margin=25mm]{geometry}
\usepackage{fontspec}
\setmainfont{Latin Modern Roman}
\usepackage{amsmath,amssymb,mathtools,unicode-math}
\setmathfont{Latin Modern Math}
\usepackage{enumitem,needspace,xurl,hyperref}
\setlength{\parindent}{0pt}
\setlength{\parskip}{5pt}
\setlength{\emergencystretch}{3em}
\newcommand{\N}{\mathbb{N}}
\newcommand{\Z}{\mathbb{Z}}
\newcommand{\Q}{\mathbb{Q}}
\newcommand{\R}{\mathbb{R}}
\newcommand{\C}{\mathbb{C}}
\newcommand{\F}{\mathbb{F}}
\newcommand{\A}{\mathbb{A}}
\newcommand{\PP}{\mathbb{P}}
\newcommand{\E}{\mathbb{E}}
\newcommand{\eps}{\varepsilon}
\newcommand{\dd}{\,\mathrm{d}}
\begin{document}
\section*{UnsolvedMath ID: 30005670}
\textbf{Source problem code:} catalogue entry
\subsection{Definitions and mathematical statement}
The supplied record contains no recoverable mathematical problem statement. Its extracted passage is a workshop abstract, so it provides no mathematical objects to define, hypotheses, domains, boundary conditions, quantifiers, or conclusion from which a precise conjecture or question can be reconstructed.

The abstract describes a workshop covering chromatic homotopy theory, algebraic \(K\)-theory, derived algebra, and equivariant homotopy theory. It reports a lecture series on the recent disproof of the telescope conjecture, 24 research talks, and two gong-shows featuring 10-minute presentations. These are contextual reports, not a formulation of the telescope conjecture or a separate open problem.

The supplied statement status is \texttt{unrecoverable}. The supplied verification field says that the official report confirms that the extraction is a workshop abstract containing no standalone open problem; this verification claim has not been independently checked here. The registry nevertheless labels the entry \texttt{open}. That label does not establish an open mathematical question, and no replacement problem is inferred.
\subsection{Short English statement}
This entry records a homotopy theory workshop abstract rather than a mathematical question. Its mention of a disproof of the telescope conjecture supplies context but no precise formulation. A statement cannot be reconstructed without additional source evidence.
\subsection{Sources}
\begin{itemize}
\item Local UnsolvedMath registry, supplied source record for permanent ID \(30005670\), source code \texttt{catalogue entry}, titled ``Equivariant Homotopy-Theory Workshop Abstract,'' in the Oberwolfach Reports Open Problems collection. The record lists proposed year 2023 and supplies no proposer attribution.
\item Supplied source citation: \emph{Homotopy Theory (2024)}, Oberwolfach Reports, DOI 10.4171/owr/2023/34, pp. 1893--1952. Supplied source URL: \url{https://doi.org/10.4171/owr/2023/34}. The citation's 2024 label and the record's 2023 year are retained without resolving their relationship.
\item Supplied literature source: ``OWR Homotopy Theory report,'' \url{https://doi.org/10.4171/OWR/2023/34}.
\item Supplied literature source: ``K-theoretic counterexamples to Ravenel's telescope conjecture,'' \url{https://arxiv.org/abs/2310.17459}. No author attribution is supplied in this record.
\item Literature assessment dated exactly 2026-08-22: the extracted text is a report abstract describing the recent disproof of the telescope conjecture, not a separate mathematical problem statement. The accompanying status is \texttt{open}. This is supplied dated curation derived from the OWR record, not an independently verified literature assessment; time-sensitive status and cited claims require independent verification.
\end{itemize}
\end{document}
